\documentclass[10pt]{amsart}
\usepackage[a4paper,margin=23mm]{geometry}
\usepackage{amsmath,amssymb,amsthm,mathtools,hyperref,graphicx,xspace,letltxmacro,xpatch}
\usepackage[english]{babel}
\usepackage{xfrac,xcolor}
\providecommand{\resp}[1]{#1}

\usepackage{mathrsfs}
\let\mathcal\mathscr

\let\epsilon\varepsilon
\newcommand{\Psfrac}[2]{(\sfrac{#1}{#2})}
\newcommand{\Pspfrac}[2]{(\sfrac{#1}{(#2)})}
\newcommand{\psfrac}[2]{\sfrac{(#1)}{#2}}
\newcommand{\spfrac}[2]{\sfrac{#1}{(#2)}}
\def\mto{\mathchoice{\longmapsto}{\mapsto}{\mapsto}{\mapsto}}

\usepackage{bbm}
\usepackage{esint}

\theoremstyle{plain}

\newtheorem{theoreme}[equation]{Theorem}
\newtheorem{proposition}[equation]{Proposition}
\newtheorem{corollary}[equation]{Corollary}
\newtheorem{lemma}[equation]{Lemma}

\theoremstyle{definition}

\newtheorem{definition}[equation]{Definition}
\newtheorem{remark}[equation]{Remark}
\newtheorem*{question}{Question}

\newcommand{\HH}{\mathbb{H}}
\newcommand{\ZZ}{\mathbb{Z}}
\newcommand{\RR}{\mathbb{R}}
\newcommand{\CC}{\mathbb{C}}
\newcommand{\NN}{\mathbb{N}}
\newcommand{\inte}[1]{\underset{#1}{\int}}
\newcommand{\finte}[3]{\underset{#1}{\fint} {#2} \, {#3}}
\newcommand{\supp}[1]{\mathrm{supp}(#1)}
\newcommand{\limiinf}[1]{\underset{#1}{\liminf}}
\newcommand{\somme}[1]{\sum_{#1}}
\newcommand{\limi}[1]{\underset{#1}{\lim}}
\newcommand{\tends}[1]{\underset{#1}{\longrightarrow}}
\newcommand{\equivaut}[1]{\underset{#1}{\sim}}
\newcommand{\fonction}[5]{\begin{array}{r r c l}
#1 \hspace{2mm} : & #2 & \displaystyle\to & #3 \\
& #4 & \displaystyle\mto & #5 \\
\end{array}}
\newcommand{\fonctionbis}[4]{\begin{split}
#1 & \to #2 \\
#3 & \mto #4
\end{split}}
\let\quotient\sfrac

\renewcommand{\resp}[1]{\textit{(resp.\ }#1\textit{)}}

\DeclareMathOperator{\vol}{vol}
\DeclareMathOperator{\Id}{Id}

\let\dpl\displaystyle
\providecommand{\skpt}{\par\smallskip\noindent}
\providecommand{\pointir}{.\ }
\providecommand{\eg}{e.g.}\providecommand{\ie}{i.e.}\providecommand{\etc}{etc.}\providecommand{\cf}{cf.}
\providecommand{\mainmatter}{}\providecommand{\backmatter}{}
\let\Subsection\subsection\let\Subsubsection\subsubsection
\AtBeginDocument{\gdef\shorttitle{Orbital functions and heat kernels}}
